\documentclass{article}
\usepackage{fontspec}
\usepackage{xunicode}
\usepackage{polyglossia}
\setmainlanguage{french}

% --- Packages ajoutés pour la correction ---
\usepackage{csquotes}       % guillemets et citations imbriquées
\usepackage{hyperref}       % hyperliens (\url, \href)
\usepackage{verse}          % environnement verse avec indentation

\begin{document}

% ============================================================
\section{csquotes}
% ============================================================

% \enquote{} pose les guillemets du niveau courant.
% Imbriqué, il passe automatiquement aux guillemets intérieurs.

Victor Hugo a écrit dans la \emph{Préface de Cromwell} :
\enquote{C'est par là qu'ils touchent à l'humanité, c'est par là
qu'ils sont dramatiques. \enquote{Du sublime au ridicule il n'y a
qu'un pas}, disait Napoléon, quand il fut convaincu d'être homme.}


% ============================================================
\section{quote}
% ============================================================

% Environnement `quote` : citation courte, un seul paragraphe,
% marges réduites des deux côtés.
%
% \textelp{}  : points de suspension entre crochets […]
% \textins{}  : insertion entre crochets [commentaire]
% Les deux commandes viennent du package csquotes.

Victor Hugo présente ainsi le public d'une pièce de théâtre :

\begin{quote}
Trois espèces de spectateurs composent ce qu'on est convenu
d'appeler le public : premièrement, les femmes ; deuxièmement,
les penseurs ; troisièmement, la foule proprement dite. Ce que
la foule demande presque exclusivement à l'œuvre dramatique,
c'est de l'action ; ce que les femmes y veulent avant tout,
c'est de la passion ; ce qu'y cherchent plus spécialement les
penseurs, ce sont des caractères.
\textelp{} \textins{c'est le rôle du dramaturge de satisfaire
ces trois publics}\footnote{Dans la \emph{Préface de Ruy Blas}.}
\end{quote}


% ============================================================
\section{quotation}
% ============================================================

% Environnement `quotation` : citation longue sur plusieurs
% paragraphes. Contrairement à `quote`, les paragraphes sont
% indentés à la première ligne.

\begin{quotation}
Homère, en effet, domine la société antique. Dans cette société,
tout est simple, tout est épique. La poésie est religion, la
religion est loi. À la virginité du premier âge a succédé la
chasteté du second. Une sorte de gravité solennelle s'est
empreinte partout, dans les mœurs domestiques comme dans les
mœurs publiques. Les peuples n'ont conservé de la vie errante
que le respect de l'étranger et du voyageur. La famille a une
patrie ; tout l'y attache ; il y a le culte du foyer, le culte
des tombeaux.

Nous le répétons, l'expression d'une pareille civilisation ne
peut être que l'épopée. L'épopée y prendra plusieurs formes,
mais ne perdra jamais son caractère. Pindare est plus sacerdotal
que patriarchal, plus épique que lyrique. Si les annalistes,
contemporains nécessaires de ce second âge du monde, se mettent
à recueillir les traditions et commencent à compter avec les
siècles, ils ont beau faire, la chronologie ne peut chasser la
poésie ; l'histoire reste épopée. Hérode est un Homère.

\raggedleft(Victor Hugo, \emph{Cromwell}, Préface)
\end{quotation}


% ============================================================
\section{url}
% ============================================================

% \url{}    : lien cliquable affiché tel quel (police machine)
% \href{}{} : lien cliquable avec texte personnalisé

Pour avoir le texte complet, suivez le lien :
\url{https://fr.wikisource.org/wiki/Cromwell_-_Préface}

% Variante avec texte personnalisé :
% \href{https://fr.wikisource.org/wiki/Cromwell_-_Préface}{Préface de Cromwell sur Wikisource}


% ============================================================
\section{verse}
% ============================================================

% Environnement `verse` (package verse) :
%   - \\ termine un vers
%   - \> indente d'une unité (pour les vers 3 et 4)
%   - bonus : \settowidth\versewidth{} centre le bloc sur le vers le plus long

Amusons-nous un peu avec Edward \textsc{Lear},
\emph{The Owl and the Pussy Cat, and Other Nonsense Poetry}.

\settowidth\versewidth{Fell casually into a kettle}
\begin{verse}[\versewidth]
There was an old man who when little \\
Fell casually into a kettle \\
\> But, growing too stout, \\
\> He could never get out, \\
So he passed all his life in that kettle
\end{verse}

\end{document}